\begin{thebibliography}{99}
\bibitem[1]{ref1} Noga Alon. Explicit expanders of every degree and size. Combinatorica, pages 1–17, 2021.
\bibitem[2]{ref2} Noga Alon, Jehoshua Bruck, Joseph Naor, Moni Naor, and Ron M. Roth. Construction of asymptotically good low-rate error-correcting codes through pseudo-random graphs. Informa- tion Theory, IEEE Transactions on, 38(2):509–516, 1992.
\bibitem[3]{ref3} Noga Alon, Oded Goldreich, Johan Håstad, and René Peralta. Simple constructions of almost k-wise independent random variables. Random Structures \& Algorithms, 3(3):289–304, 1992. New Near-Linear Time Decodable Codes Closer to the GV Bound
\bibitem[4]{ref4} Avraham Ben-Aroya and Amnon Ta-Shma. A combinatorial construction of almost-Ramanujan graphs using the zig-zag product. SIAM Journal on Computing, 40(2):267–290, 2011.
\bibitem[5]{ref5} Avraham Ben-Aroya and Amnon Ta-Shma. Constructing small-bias sets from algebraic- geometric codes. Theory of Computing, 9(5):253–272, 2013.
\bibitem[6]{ref6} Yonatan Bilu and Shlomo Hoory. On codes from hypergraphs. European Journal of Combin- atorics, 25(3):339–354, 2004.
\bibitem[7]{ref7} Yonatan Bilu and Nathan Linial. Lifts, discrepancy and nearly optimal spectral gap. Combin- atorica, 26(5):495–519, 2006.
\bibitem[8]{ref8} Andrej Bogdanov. A different way to improve the bias via expanders. Topics in (and out) the theory of computing, Lecture, 2012.
\bibitem[9]{ref9} Emma Cohen, Dhruv Mubayi, Peter Ralli, and Prasad Tetali. Inverse expander mixing for hypergraphs. The Electronic Journal of Combinatorics, 23(2):P2–20, 2016.
\bibitem[10]{ref10} David Conlon, Jonathan Tidor, and Yufei Zhao. Hypergraph expanders of all uniformities from Cayley graphs. Proceedings of the London Mathematical Society, 121(5):1311–1336, 2020.
\bibitem[11]{ref11} Joel Friedman and Avi Wigderson. On the second eigenvalue of hypergraphs. Combinatorica, 15(1):43–65, 1995.
\bibitem[12]{ref12} Edgar N. Gilbert. A comparison of signalling alphabets. The Bell system technical journal, 31(3):504–522, 1952.
\bibitem[13]{ref13} Oded Goldreich. On constructing expanders for any number of vertices, October 2019. Available at \url{https://www.wisdom.weizmann.ac.il/~oded/R1/ex4all.pdf}.
\bibitem[14]{ref14} Konstantin Golubev and Ori Parzanchevski. Spectrum and combinatorics of two-dimensional Ramanujan complexes. Israel Journal of Mathematics, 230(2):583–612, 2019.
\bibitem[15]{ref15} Venkatesan Guruswami and Piotr Indyk. Linear-time encodable/decodable codes with near- optimal rate. IEEE Transactions on Information Theory, 51(10):3393–3400, 2005.
\bibitem[16]{ref16} Venkatesan Guruswami, Atri Rudra, and Madhu Sudan. Essential Coding Theory, 2015. URL: \url{http://www.cse.buffalo.edu/faculty/atri/courses/coding-theory/book}.
\bibitem[17]{ref17} Brett Hemenway, Noga Ron-Zewi, and Mary Wootters. Local list recovery of high-rate tensor codes and applications. SIAM Journal on Computing, pages FOCS17–157, 2019.
\bibitem[18]{ref18} Wassily Hoeffding. Probability inequalities for sums of bounded random variables. Journal of the American statistical association, 58(301):13–30, 1963.
\bibitem[19]{ref19} Fernando Granha Jeronimo, Dylan Quintana, Shashank Srivastava, and Madhur Tulsiani. Unique decoding of explicit \(\varepsilon\)-balanced codes near the Gilbert–Varshamov bound. In Proceed- ings of the 61st Annual Symposium on Foundations of Computer Science (FOCS 2020), pages 434–445. IEEE, 2020.
\bibitem[20]{ref20} Fernando Granha Jeronimo, Shashank Srivastava, and Madhur Tulsiani. Near-linear time decoding of Ta-Shma’s codes via splittable regularity. In Proceedings of the 53rdth Annual Symposium on Theory of Computing (STOC 2021), pages 1527–1536. ACM, 2021.
\bibitem[21]{ref21} Swastik Kopparty, Nicolas Resch, Noga Ron-Zewi, Shubhangi Saraf, and Shashwat Silas. On list recovery of high-rate tensor codes. IEEE Transactions on Information Theory, 67(1):296–316, 2020.
\bibitem[22]{ref22} Swastik Kopparty, Noga Ron-Zewi, Shubhangi Saraf, and Mary Wootters. Improved decoding of folded Reed–Solomon and multiplicity codes. In Proceedings of the 59th Annual Symposium on Foundations of Computer Science (FOCS 2018), pages 212–223. IEEE, 2018.
\bibitem[23]{ref23} John Lenz and Dhruv Mubayi. Eigenvalues and linear quasirandom hypergraphs. In Forum of Mathematics, Sigma, volume 3. Cambridge University Press, 2015.
\bibitem[24]{ref24} Ray Li and Mary Wootters. Improved list-decodability of random linear binary codes. In APPROX–RANDOM, volume 116 of LIPIcs, pages 50:1–50:19. Schloss Dagstuhl - Leibniz- Zentrum für Informatik, 2018.
\bibitem[25]{ref25} Eyal Lubetzky and Yuval Peres. Cutoff on all Ramanujan graphs. Geometric and Functional Analysis, 26(4):1190–1216, 2016.
\bibitem[26]{ref26} Alexander Lubotzky, Ralph Phillips, and Peter Sarnak. Ramanujan graphs. Combinatorica, 8(3):261–277, 1988. 10:35
\bibitem[27]{ref27} Alexander Lubotzky, Beth Samuels, and Uzi Vishne. Explicit constructions of Ramanujan complexes of type Ad. European Journal of Combinatorics, 26(6):965–993, 2005.
\bibitem[28]{ref28} Grigorii Aleksandrovich Margulis. Explicit group-theoretical constructions of combinatorial schemes and their application to the design of expanders and concentrators. Problemy peredachi informatsii, 24(1):51–60, 1988.
\bibitem[29]{ref29} Jack Murtagh, Omer Reingold, Aaron Sidford, and Salil Vadhan. Deterministic approximation of random walks in small space. Theory of Computing, 17(1):1–35, 2021.
\bibitem[30]{ref30} Joseph Naor and Moni Naor. Small-bias probability spaces: Efficient constructions and applications. SIAM Journal on Computing, 22(4):838–856, 1993.
\bibitem[31]{ref31} Ori Parzanchevski. Mixing in high-dimensional expanders. Combinatorics, Probability and Computing, 26(5):746–761, 2017.
\bibitem[32]{ref32} Ori Parzanchevski, Ron Rosenthal, and Ran J. Tessler. Isoperimetric inequalities in simplicial complexes. Combinatorica, 36(2):195–227, 2016.
\bibitem[33]{ref33} Michael Sipser and Daniel A. Spielman. Expander codes. IEEE transactions on Information Theory, 42(6):1710–1722, 1996.
\bibitem[34]{ref34} Amnon Ta-Shma. Explicit, almost optimal, \(\varepsilon\)-balanced codes. In Proceedings of the 49th Annual Symposium on Theory of Computing (STOC 2017), pages 238–251. ACM, 2017.
\bibitem[35]{ref35} R. Tanner. A recursive approach to low complexity codes. IEEE Transactions on information theory, 27(5):533–547, 1981.
\bibitem[36]{ref36} Rom Rubenovich Varshamov. Estimate of the number of signals in error correcting codes. Docklady Akad. Nauk, SSSR, 117:739–741, 1957.
\bibitem[37]{ref37} Gillés Zémor. On expander codes. IEEE Transactions on Information Theory, 47(2):835–837, 2001.
\end{thebibliography}
